\documentclass[letterpaper, 10pt]{article}

\usepackage{luacode}

\usepackage[left=1in, top=1in, total={6.5in, 9in}]{geometry}

\usepackage{booktabs}

\usepackage{longtable}

\usepackage{hyperref}
\hypersetup{colorlinks, urlcolor=blue}
\urlstyle{same}

\usepackage{array}

\usepackage[]{ltablex}
\usepackage{makecell}
\renewcommand{\cellalign}{tl}

\setlength\LTleft{0pt}
\setlength\LTright{0pt}
\setlength{\LTpre}{0pt}
\setlength{\LTpost}{0pt}

\usepackage{titlesec}

\titleformat*{\section}{\large\bfseries}

\newcommand{\gutterSpacing}{0.125\textwidth}

\newcolumntype{D}{>{\itshape\arraybackslash}p{\gutterSpacing}}
\newcolumntype{R}{>{}p{\gutterSpacing}}

\usepackage{hyphenat}

\usepackage{lastpage}
\usepackage{fancyhdr}
\pagestyle{fancy}
\fancyhf{}


\usepackage{siunitx}

% Custom command

\newcommand{\headings}[1]{\section*{#1} \hrule \vspace{10pt}}
\newcommand{\subheadings}[1]{\textbf{#1} \vspace{5pt}}

\begin{document}

% Ensures that the first page doesn't have the header...
\thispagestyle{empty}

\begin{luacode}
	-- This builds the FULL CV by design: every entry in data/*.json is printed.
	-- The "featured" flag present in some entries is used only by the Hugo
	-- website (layouts/) to decide what to highlight on the homepage; it is
	-- intentionally ignored here.

	-- FUNCTIONS
	-- A nice little function to test if a string is empty
	local function isEmpty(s)
		return s == nil or s == ''
	end
	-- A little function to see if an array contains a value
	local function matches(thisValue, arrayToMatch)
		for index, value in pairs(arrayToMatch) do
			if value == thisValue then
				return true
			end
		end
		return false
	end
	-- A little function to capitalize the first letter in a string
	function firstToUpper(str)
		return (str:gsub("^%l", string.upper))
	end

	-- For human-readable text (never for the URL argument of an href link)
	function prepareString(str)
		if isEmpty(str) then return '' end
		return (str:gsub('([&%%#_])', '\\%1'))
	end
	-- For the URL argument of an href link only. URLs printed from Lua via
	-- tex.print are tokenized with normal catcodes, so ampersand, percent and
	-- hash must be backslash-escaped; hyperref maps the escaped forms back
	-- to the literal characters in the target URL.
	function prepareUrl(str)
		if isEmpty(str) then return '' end
		return (str:gsub('([&%%#])', '\\%1'))
	end

	function getYear(str)
		if isEmpty(str) then return nil end
		local pattern = "(%d+)-(%d+)-(%d+)"
		local year, month, day = str:match(pattern)
		return year
	end

	function getMonth(str)
		if isEmpty(str) then return nil end
		local pattern = "(%d+)-(%d+)-(%d+)"
		local year, month, day = str:match(pattern)
		return month
	end

	function getYearAndMonth(str)
		if isEmpty(str) then return nil end
		local pattern = "(%d+)-(%d+)-(%d+)"
		local year, month, day = str:match(pattern)
		if year == nil then return nil end
		return year .. '-' .. month
	end

	-- fmt is getYear or getYearAndMonth
	function formatRange(startStr, endStr, fmt)
		local s, e = fmt(startStr), fmt(endStr)
		if s == nil then return e or '' end
		if e == nil then return s .. ' - Present' end
		if s == e then return s end
		return s .. ' - ' .. e
	end

	function replaceEmpty(str)
		if isEmpty(str) then
			return ''
		else
			return str
		end
	end

	-- Nil-tolerant: entries with a missing date sort as the smallest (last
	-- in the descending sort order used throughout this document).
	function compareYear(a, b, flag)
		local av, bv = a[flag], b[flag]
		if isEmpty(av) then
			return false
		elseif isEmpty(bv) then
			return true
		else
			return av > bv
		end
	end

	function compareYearYear(a, b)
		return compareYear(a, b, "year")
	end

	function compareYearStart(a, b)
		return compareYear(a, b, "start")
	end

	function compareYearMonth(a, b, flag)
		local ay, by = getYear(a[flag]), getYear(b[flag])
		if isEmpty(ay) and isEmpty(by) then
			return false
		elseif isEmpty(ay) then
			return false
		elseif isEmpty(by) then
			return true
		elseif not (ay == by) then
			return ay > by
		else
			local am, bm = getMonth(a[flag]), getMonth(b[flag])
			if isEmpty(am) then
				return false
			elseif isEmpty(bm) then
				return true
			else
				return am > bm
			end
		end
	end

	function compareYearMonthGraduation(a, b)
		return compareYearMonth(a, b, "graduation")
	end

	function compareYearMonthEnd(a, b)
		return compareYearMonth(a, b, "end")
	end

	function compareYearMonthStart(a, b)
		return compareYearMonth(a, b, "start")
	end
	-- END FUNCTIONS

	-- METADATA
	require("lualibs.lua")

	local function readJson(path)
		local f = io.open(path, 'r')
		assert(f, path .. ' not found; run from the repository root')
		local raw = f:read('*a')
		f:close()
		return utilities.json.tolua(raw)
	end

	local function readJsonOptional(path)
		local f = io.open(path, 'r')
		if not f then return nil end
		local raw = f:read('*a')
		f:close()
		return utilities.json.tolua(raw)
	end

	local my = readJson('data/mydata.json')

	-- HEADER
	tex.print{'\\lhead{' .. prepareString(my["name"]) .. '}'}
	tex.print{'\\rhead{\\textit{Curriculum Vitae, \\thepage\\ of \\pageref*{LastPage}}}'}
	tex.print{'\\cfoot{}'}

	tex.print('\\begin{center}')
	tex.print('{\\Large \\textbf{' .. prepareString(my["name"]) .. '}}\\\\')
	tex.print(prepareString(my["department"]) .. '\\\\')
	tex.print(prepareString(my["university"]) .. '\\\\')
	if not isEmpty(my["address1"]) then
		tex.print(prepareString(my["address1"]) .. '\\\\')
	end
	if not isEmpty(my["address2"]) then
		tex.print(prepareString(my["address2"]) .. '\\\\')
	end
	if not isEmpty(my["address3"]) then
		tex.print(prepareString(my["address3"]) .. '\\\\')
	end
	-- Optional: a site that publishes a contact form instead of an address
	-- omits "email" from mydata.json, and the CV simply carries the website.
	if not isEmpty(my["email"]) then
		tex.print('\\href{mailto:' .. prepareUrl(my["email"]) .. '}{' .. prepareString(my["email"]) .. '} \\\\')
	end
	tex.print('\\href{' .. prepareUrl(my["website"]) .. '}{' .. prepareString(my["website"]) .. '} \\\\')
	tex.print('\\end{center}')

	-- My name (for citations)
	local myName = my["citationName"]
	local myNamePattern = myName:gsub('([%.%-%%])', '%%%1')
	local myNameStylized = '\\textbf{' .. myName .. '}'


	-- local prettyprintPublications
	prettyprintPublications = function (data, section)
		tex.print('\\headings{' .. prepareString(section.cvLabel) .. '}')

		-- Publications are challenging to lay out
		-- We begin by assembling tables
		-- Add each manuscript status you want to show, in the order that you want it to come out
		local orderedStates = {'thesis', 'unpublished', 'article', 'incollection', 'inproceeding'}

		-- Count only the entries that will actually be printed (i.e. whose
		-- entry-type is one of orderedStates) so the descending numbering
		-- below is contiguous.
		local pubCount = 0
		for dataKey, dataValue in ipairs(data) do
			if matches(dataValue["entry-type"], orderedStates) then
				pubCount = pubCount + 1
			end
		end

		-- Sort the table by year
		function compare(a, b)
			if isEmpty(a["description"]["year"]) then
				return false
			elseif isEmpty(b["description"]["year"]) then
				return true
			else
				return a["description"]["year"] > b["description"]["year"]
			end
		end

		for _, thisOrderedState in ipairs(orderedStates) do 											-- loop through all types
			local needsHeading = true

			-- make table for thisOrderedState
			local stateData = {}
			for dataKey, dataValue in ipairs(data) do
				if dataValue["entry-type"] == thisOrderedState then
					table.insert(stateData, dataValue)
				end
			end

			-- sort by year
			table.sort(stateData, compare)

			-- loop over state data
			for pubIdx, pubState in ipairs(stateData) do
				if needsHeading then
					tex.print('\\begin{flushleft}')
					if (pubState["entry-type"] == 'thesis') then
						tex.print("\\subheadings{Theses}")
					elseif (pubState["entry-type"] == 'unpublished') then
						tex.print("\\subheadings{Preprints}")
					elseif (pubState["entry-type"] == 'article') then
						tex.print("\\subheadings{Journal articles}")
					elseif (pubState["entry-type"] == 'incollection') then
						tex.print("\\subheadings{Book chapters}")
					elseif (pubState["entry-type"] == 'inproceeding') then
						tex.print("\\subheadings{Conference papers}")
					else
						tex.print("\\subheadings{" .. firstToUpper(pubState["entry-type"]) .. "}")
					end
					tex.print('\\end{flushleft}')
					tex.print('\\begin{centering}')
					tex.print('\\renewcommand{\\arraystretch}{1.5}')
					tex.print('\\begin{longtable}{p{0.1\\textwidth}p{0.9\\textwidth}}')

					needsHeading = false
				end

				local thisAuthor 	= prepareString(pubState["description"]["author"])
				local thisTitle 	= prepareString(pubState["description"]["title"])

				assert(not isEmpty(thisAuthor),"Author must be provided")
				assert(not isEmpty(thisTitle),"Title must be provided")

				-- URL is optional
				local thisUrl = pubState["description"]["url"]
				if isEmpty(thisUrl) then
					thisUrl = ''
				else
					thisUrl = '(\\href{' .. prepareUrl(thisUrl) .. '}{link})'
				end

				local thisCitation = ''
				-- THESIS SPECIALIZED
				if pubState['entry-type'] == 'thesis' then
					local thisSchool = prepareString(pubState["description"]["school"])
					local thisDate = prepareString(pubState["description"]["date"])
					local thisDegree = prepareString(pubState["description"]["degree"])

					assert(not isEmpty(thisSchool),"School must be provided")
					assert(not isEmpty(thisDate),"Date must be provided")
					assert(not isEmpty(thisDegree),"Degree must be provided")

					-- FORMAT CITATION
					thisCitation = thisAuthor .. ', ``' .. thisTitle .. "''" ..
					'\\newline \\textit{' .. thisDegree .. '}' .. ' (' .. thisDate .. ') ' .. thisSchool ..
					'.\\newline' .. thisUrl


				-- PREPRINT SPECIALIZED
				elseif pubState['entry-type'] == 'unpublished' then
					local thisEprint = prepareString(replaceEmpty(pubState["description"]["eprint"]))
					if isEmpty(thisEprint) then
						thisEprint = ''
					else
						thisEprint = '. \\textit{' .. thisEprint .. '}'
					end

					-- FORMAT CITATION
					thisCitation = thisAuthor .. ', ``' .. thisTitle .. "''" .. thisEprint ..
					'.\\newline' .. thisUrl

				elseif pubState['entry-type'] == 'article' then
					local thisYear = prepareString(pubState["description"]["year"])
					local thisJournal = prepareString(pubState["description"]["journal"])

					assert(not isEmpty(thisYear),"Year must be provided")
					assert(not isEmpty(thisJournal),"Journal must be provided")

					-- Pages and volume are optional
					local thisPages = prepareString(replaceEmpty(pubState["description"]["pages"]))
					local thisVolume = prepareString(replaceEmpty(pubState["description"]["volume"]))

					local thisVolumePart = ''
					if not isEmpty(thisVolume) then
						thisVolumePart = ', ' .. thisVolume
					end
					local thisPagesPart = ''
					if not isEmpty(thisPages) then
						thisPagesPart = ', pp. ' .. thisPages
					end

					-- FORMAT CITATION
					thisCitation = thisAuthor .. ' (' .. thisYear .. '). ' .. thisTitle ..
					'. \\textit{' .. thisJournal .. '}' .. thisVolumePart .. thisPagesPart .. '.' ..
					'\\newline' .. thisUrl


				elseif pubState['entry-type'] == 'inproceeding' or pubState['entry-type'] == 'incollection' then
					local thisYear = prepareString(pubState["description"]["year"])

					assert(not isEmpty(thisYear),"Year must be provided")

					-- Pages, editor and publisher are optional; only the
					-- booktitle is required.
					local thisPages = prepareString(replaceEmpty(pubState["description"]["pages"]))
					local thisBooktitle = prepareString(replaceEmpty(pubState["description"]["booktitle"]))
					local thisPublisher = prepareString(replaceEmpty(pubState["description"]["publisher"]))
					local thisEditor = prepareString(replaceEmpty(pubState["description"]["editor"]))

					assert(not isEmpty(thisBooktitle),"Booktitle must be provided")

					local thisEditorPart = ''
					if not isEmpty(thisEditor) then
						thisEditorPart = thisEditor .. ' (Ed.), '
					end

					local thisPagesPart = ''
					if not isEmpty(thisPages) then
						thisPagesPart = ' (pp. ' .. thisPages .. ')'
					end

					local thisPublisherPart = ''
					if not isEmpty(thisPublisher) then
						thisPublisherPart = '. ' .. thisPublisher
					end

					-- FORMAT CITATION
					thisCitation = thisAuthor .. ' (' .. thisYear .. '). ' .. thisTitle ..
					' In ' .. thisEditorPart .. thisBooktitle ..
					thisPagesPart .. thisPublisherPart .. '.' ..
					'\\newline' .. thisUrl

				end

				-- let's make my name bold
				thisCitation = thisCitation:gsub(myNamePattern, myNameStylized)

				tex.print(pubCount .. '. & ' .. thisCitation .. '\\\\')

				pubCount = pubCount - 1
			end

			-- Only close the table if one was opened; see above.
			if not needsHeading then
				tex.print('\\end{longtable}')
				tex.print('\\end{centering}')
			end
		end
	end

	-- local prettyprintEducation
	prettyprintEducation = function (data, section)
		tex.print('\\headings{' .. prepareString(section.cvLabel) .. '}')

		tex.print('\\begin{center}')
		tex.print('\\renewcommand{\\arraystretch}{1.5}')
		tex.print('\\begin{longtable}{p{0.2\\textwidth}p{0.8\\textwidth}}')
		for dataKey, dataValue in ipairs(data) do
			local thisStartdate = dataValue["start"]
			local thisEnddate   = dataValue["end"]

			local columnOne = formatRange(thisStartdate, thisEnddate, getYear)

			local thisSchool    	= prepareString(replaceEmpty(dataValue["school"]))
			local thisDegree    	= prepareString(replaceEmpty(dataValue["degree"]))
			local thisMajor    		= prepareString(replaceEmpty(dataValue["major"]))
			local thisThesis 		= prepareString(replaceEmpty(dataValue["thesis"]))
			local thisNote 			= prepareString(replaceEmpty(dataValue["notes"]))
			local thisSupervisors 	= prepareString(replaceEmpty(dataValue["supervisors"]))
			local thisLink 			= replaceEmpty(dataValue["thesislink"])

			assert(not isEmpty(thisDegree),"A degree is missing")
			assert(not isEmpty(thisSchool),"A school is missing")

			thisSchool = '\\textsc{' .. thisSchool .. '}'

			if not isEmpty(thisMajor) then
				thisMajor = ' in ' .. thisMajor
			end
			if not isEmpty(thisNote) then
				thisNote = '(' .. thisNote .. ')'
			end
			if not (isEmpty(thisThesis) or isEmpty(thisSupervisors) or isEmpty(thisLink)) then
				thisThesis = '\\newline ' .. "\\textit{Thesis}: ``" .. thisThesis .. "''" .. '\\newline '
				thisSupervisors = "\\textit{Supervisors}: " .. thisSupervisors
			end

			local thisContents = thisDegree .. thisMajor .. ' ' .. thisNote .. '\\newline ' ..
								 thisSchool ..
								 thisThesis ..
								 thisSupervisors

			tex.print(columnOne .. ' & ' .. thisContents .. ' \\\\')
		end
		tex.print('\\end{longtable}')
		tex.print('\\end{center}')
	end

	-- local prettyprintExperience
	prettyprintExperience = function (data, section)
		tex.print('\\headings{' .. prepareString(section.cvLabel) .. '}')

		-- Add each manuscript status you want to show, in the order that you want it to come out
		local orderedStates = {'professional', 'visit', 'voluntary'}

		for _, thisOrderedState in ipairs(orderedStates) do 											-- loop through all types
			local needsHeading = true

			-- make table for thisOrderedState
			local stateData = {}
			for dataKey, dataValue in ipairs(data) do
				if dataValue["type"] == thisOrderedState then
					table.insert(stateData, dataValue)
				end
			end

			-- sort by year
			table.sort(stateData, compareYearMonthStart)

			-- loop over state data
			for expIdx, expState in ipairs(stateData) do
				if needsHeading then
					tex.print('\\begin{flushleft}')
					if (expState["type"] == 'professional') then
						tex.print("\\subheadings{Professional positions}")
					elseif (expState["type"] == 'visit') then
						tex.print("\\subheadings{Research visits}")
					elseif (expState["type"] == 'voluntary') then
						tex.print("\\subheadings{Voluntary work}")
					else
						tex.print("\\subheadings{" .. firstToUpper(expState["type"]) .. "}")
					end
					tex.print('\\end{flushleft}')
					tex.print('\\begin{centering}')
					tex.print('\\renewcommand{\\arraystretch}{1.5}')
					tex.print('\\begin{longtable}{p{0.2\\textwidth}p{0.8\\textwidth}}')
					needsHeading = false
				end

				local columnOne = formatRange(expState["start"], expState["end"], getYearAndMonth)

				if expState["type"] == 'professional' then
					local thisInstitute 	= prepareString(replaceEmpty(expState["institute"]))
					local thisDepartment    = prepareString(replaceEmpty(expState["department"]))
					local thisLocation      = prepareString(replaceEmpty(expState["location"]))
					local thisRole   		= prepareString(replaceEmpty(expState["role"]))

					assert(not isEmpty(thisRole),"A role is missing")
					assert(not isEmpty(thisInstitute),"An institute is missing")
					assert(not isEmpty(thisDepartment),"A department is missing")
					assert(not isEmpty(thisLocation),"A location is missing")

					thisInstitute = '\\textsc{' .. thisInstitute .. '}'

					local thisContents = thisRole ..
										'\\newline ' .. thisInstitute ..
										'\\newline ' .. thisDepartment ..
										'\\newline ' .. thisLocation

					tex.print(columnOne .. ' & ' .. thisContents .. ' \\\\')
				elseif expState["type"] == 'visit' then
					local thisInstitute 	= prepareString(replaceEmpty(expState["institute"]))
					local thisDepartment    = prepareString(replaceEmpty(expState["department"]))
					local thisLocation      = prepareString(replaceEmpty(expState["location"]))
					local thisRole   		= prepareString(replaceEmpty(expState["role"]))
					local thisVisiting    	= prepareString(replaceEmpty(expState["visiting"]))
					local thisDuration    	= prepareString(replaceEmpty(expState["duration"]))

					assert(not isEmpty(thisRole),"A role is missing")
					if isEmpty(thisInstitute) then
						thisInstitute = ''
					else
						thisInstitute = '\\newline \\textsc{' .. thisInstitute .. '}'
					end
					if isEmpty(thisDepartment) then
						thisDepartment = ''
					else
						thisDepartment = ', ' .. thisDepartment
					end
					if isEmpty(thisLocation) then
						thisLocation = ''
					else
						thisLocation = '\\newline ' .. thisLocation
					end
					if isEmpty(thisVisiting) then
						thisVisiting = ''
					else
						thisVisiting = '\\newline\\textit{Visiting}: ' .. thisVisiting
					end
					if isEmpty(thisDuration) then
						thisDuration = ''
					else
						thisDuration = ' (' .. thisDuration .. ')'
					end

					local thisContents = thisRole ..
										thisInstitute .. thisDepartment .. thisDuration ..
										thisLocation ..
										thisVisiting

					tex.print(columnOne .. ' & ' .. thisContents .. ' \\\\')
				elseif expState["type"] == 'voluntary' then
					local thisOrganization 	= prepareString(replaceEmpty(expState["organization"]))
					local thisRole  		= prepareString(replaceEmpty(expState["role"]))
					local thisLocation      = prepareString(replaceEmpty(expState["location"]))
					local thisUrl 			= replaceEmpty(expState["url"])

					assert(not isEmpty(thisOrganization),"An organization is missing")
					assert(not isEmpty(thisRole),"A role is missing")

					thisOrganization = '\\textsc{' .. thisOrganization .. '}'
					if isEmpty(thisUrl) then
						thisUrl = ''
					else
						thisUrl = '(\\href{' .. prepareUrl(thisUrl) .. '}{link})'
					end
					if isEmpty(thisLocation) then
						thisLocation = ''
					else
						thisLocation = '\\newline ' .. thisLocation
					end

					local thisContents = thisRole ..
										'\\newline ' .. thisOrganization ..
										thisLocation

					tex.print(columnOne .. ' & ' .. thisContents .. '\\\\')
				end
			end
			-- Only close the table if one was opened; see above.
			if not needsHeading then
				tex.print('\\end{longtable}')
				tex.print('\\end{centering}')
			end
		end
	end

	-- local prettyprintConferences
	prettyprintConferences = function (data, section)
		tex.print('\\headings{' .. prepareString(section.cvLabel) .. '}')

		-- make state vector for invited
		local invitedStates = {true, false}

		for _, thisInvitedState in ipairs(invitedStates) do 											-- loop through all types
			local needsHeading = true

			-- make table for thisInvitedState
			local stateData = {}
			for dataKey, dataValue in ipairs(data) do
				if dataValue["invited"] == thisInvitedState then
					table.insert(stateData, dataValue)
				end
			end

			-- sort by year
			table.sort(stateData, compareYearMonthStart)

			-- loop over state data
			for confIdx, confState in ipairs(stateData) do
				if needsHeading then
					tex.print('\\begin{flushleft}')
					if (confState["invited"]) then
						tex.print("\\subheadings{Invited presentations}")
					else
						tex.print("\\subheadings{{Contributed presentations}}")
					end
					tex.print('\\end{flushleft}')
					tex.print('\\begin{centering}')
					tex.print('\\renewcommand{\\arraystretch}{1.5}')
					tex.print('\\begin{longtable}{p{0.25\\textwidth}p{0.7\\textwidth}}')
					needsHeading = false
				end

				local thisTitle     = prepareString(confState["title"])
				local thisName      = prepareString(confState["name"])
				local thisLocation  = prepareString(confState["location"])

				assert(not isEmpty(thisTitle),"A title is missing")
				assert(not isEmpty(thisName),"A name is missing")
				assert(not isEmpty(thisLocation),"A location is missing")

				-- Conference talks are dated events, not open-ended ranges: an
				-- empty end date just means a single-day talk, never "Present",
				-- so this deliberately does not reuse formatRange's fallback.
				local thisConfStart = getYearAndMonth(confState["start"])
				local thisConfEnd = getYearAndMonth(confState["end"])
				local columnOne = thisConfStart
				if not isEmpty(thisConfEnd) then
					if not (thisConfEnd == thisConfStart) then
						columnOne = columnOne .. ' - ' .. thisConfEnd
					end
				end

				local thisContents = thisTitle ..
									'\\newline ' .. thisName ..
									'\\newline ' .. thisLocation

				tex.print(columnOne .. ' & ' .. thisContents .. ' \\\\')
			end
			-- Only close the table if one was opened. A group with no entries
			-- never opens a longtable, and closing one that was never opened
			-- aborts the run -- so a CV whose publications happen not to span
			-- every entry type would fail to build at all. (No control
			-- sequences in this comment: luacode scans the body textually for
			-- its terminator, so a stray one here ends the environment early.)
			if not needsHeading then
				tex.print('\\end{longtable}')
				tex.print('\\end{centering}')
			end
		end
	end

	-- local prettyprintTeaching
	prettyprintTeaching = function (data, section)
		tex.print('\\headings{' .. prepareString(section.cvLabel) .. '}')

		-- make state vector
		local teachingStates = {"MSc","BSc"}

		for _, thisTeachingState in ipairs(teachingStates) do
			local needsHeading = true

			-- make table for thisTeachingState
			local stateData = {}
			for dataKey, dataValue in ipairs(data) do
				if dataValue["level"] == thisTeachingState then
					table.insert(stateData, dataValue)
				end
			end

			-- sort by year
			table.sort(stateData, compareYearMonthEnd)

			-- loop over state data
			for teachIdx, teachState in ipairs(stateData) do
				if needsHeading then
					tex.print('\\begin{flushleft}')
					if (teachState["level"] == 'MSc') then
						tex.print("\\subheadings{MSc level}")
					elseif (teachState["level"] == 'BSc') then
						tex.print("\\subheadings{BSc level}")
					else
						tex.print('\\subheadings{'..firstToUpper(teachState["level"])..'}')
					end
					tex.print('\\end{flushleft}')
					tex.print('\\begin{centering}')
					tex.print('\\renewcommand{\\arraystretch}{1.5}')
					tex.print('\\begin{longtable}{p{0.2\\textwidth}p{0.8\\textwidth}}')
					needsHeading = false
				end

				local thisStartdate = teachState["start"]
				local thisEnddate   = teachState["end"]
				assert(not isEmpty(thisStartdate),"A start date is missing")

				local thisUniversity = prepareString(replaceEmpty(teachState["university"]))
				local thisDegree	 = prepareString(replaceEmpty(teachState["degree"]))
				local thisLevel 	= replaceEmpty(teachState["level"])
				local thisCourse 	= prepareString(replaceEmpty(teachState["course"]))
				local thisDuration 	= prepareString(replaceEmpty(teachState["duration"]))
				local thisRole 		= prepareString(replaceEmpty(teachState["role"]))
				local thisUrl 		= replaceEmpty(teachState["url"])

				assert(not isEmpty(thisUniversity),"A university is missing")
				assert(not isEmpty(thisDegree),"A degree is missing")
				assert(not isEmpty(thisLevel),"A level is missing")
				assert(not isEmpty(thisCourse),"A course is missing")
				assert(not isEmpty(thisDuration),"A duration is missing")
				assert(not isEmpty(thisRole),"A role is missing")

				if isEmpty(thisUrl) then
					thisUrl = ''
				else
					thisUrl = '(\\href{' .. prepareUrl(thisUrl) .. '}{link})'
				end

				local columnOne = formatRange(thisStartdate, thisEnddate, getYearAndMonth)

				local thisContents = thisCourse ..
									'\\newline\\textsc{' .. thisUniversity .. '}' ..
									'\\newline ' .. thisDegree ..
									'\\newline\\textit{Duration:} ' .. thisDuration ..
									'\\newline\\textit{Role:} ' .. thisRole ..
									'\\newline ' .. thisUrl

				tex.print(columnOne .. ' & ' .. thisContents .. ' \\\\')
			end
			-- Only close the table if one was opened. A group with no entries
			-- never opens a longtable, and closing one that was never opened
			-- aborts the run -- so a CV whose publications happen not to span
			-- every entry type would fail to build at all. (No control
			-- sequences in this comment: luacode scans the body textually for
			-- its terminator, so a stray one here ends the environment early.)
			if not needsHeading then
				tex.print('\\end{longtable}')
				tex.print('\\end{centering}')
			end
		end
	end

	-- local prettyprintAwards
	prettyprintAwards = function (data, section)
		tex.print('\\headings{' .. prepareString(section.cvLabel) .. '}')

		-- make state vector
		local awardState = {"award","membership"}

		for _, thisAwardState in ipairs(awardState) do 											-- loop through all types
			local needsHeading = true
			-- make table for thisAwardState
			local stateData = {}
			for dataKey, dataValue in ipairs(data) do
				if dataValue["type"] == thisAwardState then
					table.insert(stateData, dataValue)
				end
			end

			-- sort by year
			if thisAwardState == "award" then
				table.sort(stateData, compareYearYear)
			elseif thisAwardState == "membership" then
				table.sort(stateData, compareYearStart)
			else
				table.sort(stateData, compareYearMonthEnd)
			end

			-- loop over state data
			for awardIdx, awardState in ipairs(stateData) do
				if needsHeading then
					tex.print('\\begin{flushleft}')
					if (awardState["type"] == 'award') then
						tex.print("\\subheadings{Awards}")
					elseif (awardState["type"] == 'membership') then
						tex.print("\\subheadings{Memberships}")
					else
						tex.print('\\subheadings{'..firstToUpper(awardState["type"])..'}')
					end
					tex.print('\\end{flushleft}')
					tex.print('\\begin{centering}')
					tex.print('\\renewcommand{\\arraystretch}{1.5}')
					tex.print('\\begin{longtable}{p{0.2\\textwidth}p{0.8\\textwidth}}')
					needsHeading = false
				end

				if awardState["type"]=="award" then
					local thisTitle    	= prepareString(replaceEmpty(awardState["title"]))
					local thisLocation  = prepareString(replaceEmpty(awardState["location"]))
					local thisYear   	= replaceEmpty(awardState["year"])

					assert(not isEmpty(thisTitle),"A title is missing")
					assert(not isEmpty(thisLocation),"A location is missing")
					assert(not isEmpty(thisYear),"A year is missing")

					local columnOne = thisYear
					local thisContents = thisTitle ..
										'\\newline ' .. thisLocation

					tex.print(columnOne .. ' & ' .. thisContents .. ' \\\\')
				elseif awardState["type"]=="membership" then
					local thisOrganization 	= prepareString(replaceEmpty(awardState["organization"]))
					local thisMember   		= prepareString(replaceEmpty(awardState["member"]))
					local thisUrl 			= replaceEmpty(awardState["url"])

					assert(not isEmpty(thisOrganization),"An organization is missing")
					assert(not isEmpty(thisMember),"A member is missing")

					if isEmpty(thisUrl) then
						thisUrl = ''
					else
						thisUrl = '(\\href{' .. prepareUrl(thisUrl) .. '}{link})'
					end

					local columnOne = formatRange(awardState["start"], awardState["end"], getYear)
					local thisContents = thisMember ..
										'\\newline ' .. thisOrganization ..
										'\\newline ' .. thisUrl

					tex.print(columnOne .. ' & ' .. thisContents .. ' \\\\')
				else
					error("Unknown award type " .. awardState["type"])
				end
			end
			-- Only close the table if one was opened; see above.
			if not needsHeading then
				tex.print('\\end{longtable}')
				tex.print('\\end{centering}')
			end
		end
	end

	-- local prettyprintSoftware
	prettyprintSoftware = function (data, section)
		tex.print('\\headings{' .. prepareString(section.cvLabel) .. '}')

		tex.print('\\begin{centering}')
		tex.print('\\renewcommand{\\arraystretch}{1.5}')
		tex.print('\\begin{longtable}{p{\\textwidth}}')

		for softwareIdx, softwareState in ipairs(data) do
			local thisName		= prepareString(replaceEmpty(softwareState["name"]))
			local thisRole   	= prepareString(replaceEmpty(softwareState["role"]))
			local thisUrl    	= replaceEmpty(softwareState["url"])
			local thisDescription= prepareString(replaceEmpty(softwareState["description"]))

			assert(not isEmpty(thisName),"A name is missing")
			assert(not isEmpty(thisRole),"A role is missing")

			if (isEmpty(thisUrl)) then
				thisUrl = ''
			else
				thisUrl = '(\\href{' .. prepareUrl(thisUrl) .. '}{link})'
			end

			if isEmpty(thisDescription) then
				thisDescription = ''
			else
				thisDescription = '\\newline\\textit{Description: }' .. thisDescription
			end

			local thisContents = thisName .. ' ' .. thisUrl ..
								'\\newline\\textit{Role: }' .. thisRole ..
								thisDescription

			tex.print(thisContents .. ' \\\\')
		end
		tex.print('\\end{longtable}')
		tex.print('\\end{centering}')
	end

	-- The fallback renderer, mirroring layouts/partials/sections/generic.html:
	-- any array of flat records, laid out from the field names declared in the
	-- section's registry entry. A section invented in data/sections.json shows
	-- up here with no LaTeX written for it.
	prettyprintGeneric = function (data, section)
		tex.print('\\headings{' .. prepareString(section.cvLabel) .. '}')

		local titleField    = section.titleField or 'title'
		local subtitleField = section.subtitleField or 'subtitle'
		local bodyField     = section.bodyField or 'description'
		local dateField     = section.dateField or 'date'
		local linkField     = section.linkField or 'url'

		tex.print('\\begin{centering}')
		tex.print('\\renewcommand{\\arraystretch}{1.5}')
		tex.print('\\begin{longtable}{p{0.2\\textwidth}p{0.8\\textwidth}}')

		for _, record in ipairs(data) do
			local thisTitle    = prepareString(replaceEmpty(record[titleField]))
			local thisSubtitle = prepareString(replaceEmpty(record[subtitleField]))
			local thisBody     = prepareString(replaceEmpty(record[bodyField]))
			local thisUrl      = replaceEmpty(record[linkField])
			local thisDate     = replaceEmpty(record[dateField])

			-- Dates are ISO strings in some files and bare years in others;
			-- getYear returns nil for the latter, so keep the raw value.
			local columnOne = getYear(thisDate) or thisDate

			if isEmpty(thisUrl) then
				thisUrl = ''
			else
				thisUrl = ' (\\href{' .. prepareUrl(thisUrl) .. '}{link})'
			end
			if not isEmpty(thisSubtitle) then
				thisSubtitle = '\\newline\\textsc{' .. thisSubtitle .. '}'
			end
			if not isEmpty(thisBody) then
				-- The space matters: without it TeX reads \newline and the
				-- first word of the body as one undefined control sequence.
				thisBody = '\\newline ' .. thisBody
			end

			tex.print(columnOne .. ' & ' .. thisTitle .. thisUrl .. thisSubtitle .. thisBody .. ' \\\\')
		end
		tex.print('\\end{longtable}')
		tex.print('\\end{centering}')
	end

	-- local prettyprintSupervision
	prettyprintSupervision = function (data, section)
		tex.print('\\headings{' .. prepareString(section.cvLabel) .. '}')

		-- make state vector
		local supervisionStates = {"PhD", "MSc", "BSc"}

		for _, thisSupervisionState in ipairs(supervisionStates) do 											-- loop through all types
			local needsHeading = true

			-- make table for thisSupervisionState
			local stateData = {}
			for dataKey, dataValue in ipairs(data) do
				if dataValue["level"] == thisSupervisionState then
					table.insert(stateData, dataValue)
				end
			end

			-- sort by year
			table.sort(stateData, compareYearMonthGraduation)

			-- loop over state data
			for supervisionIdx, supervisionState in ipairs(stateData) do
				if needsHeading then
					tex.print('\\begin{flushleft}')
					if (supervisionState["level"] == 'PhD') then
						tex.print("\\subheadings{PhD level}")
					elseif (supervisionState["level"] == 'MSc') then
						tex.print("\\subheadings{MSc level}")
					elseif (supervisionState["level"] == 'BSc') then
						tex.print("\\subheadings{BSc level}")
					else
						tex.print('\\subheadings{'..firstToUpper(supervisionState["level"])..'}')
					end
					tex.print('\\end{flushleft}')
					tex.print('\\begin{centering}')
					tex.print('\\renewcommand{\\arraystretch}{1.5}')
					tex.print('\\begin{longtable}{p{0.2\\textwidth}p{0.8\\textwidth}}')
					needsHeading = false
				end

				local thisGraduation = getYearAndMonth(supervisionState["graduation"])
				assert(not isEmpty(thisGraduation),"A graduation date is missing")

				local thisName		= prepareString(replaceEmpty(supervisionState["name"]))
				local thisUniversity = prepareString(replaceEmpty(supervisionState["university"]))
				local thisDegree	 = prepareString(replaceEmpty(supervisionState["degree"]))
				local thisTitle 	 = prepareString(replaceEmpty(supervisionState["title"]))
				local thisRole   	 = prepareString(replaceEmpty(supervisionState["role"]))
				local thisCosupervisors = prepareString(replaceEmpty(supervisionState["cosupervisors"]))
				local thisUrl 		 = replaceEmpty(supervisionState["url"])

				assert(not isEmpty(thisName),"A name is missing")
				assert(not isEmpty(thisUniversity),"A university is missing")
				assert(not isEmpty(thisDegree),"A degree is missing")
				assert(not isEmpty(thisTitle),"A title is missing")
				assert(not isEmpty(thisRole),"A role is missing")

				if isEmpty(thisUrl) then
					thisUrl = ''
				else
					thisUrl = '(\\href{' .. prepareUrl(thisUrl) .. '}{link})'
				end

				if isEmpty(thisCosupervisors) then
					thisCosupervisors = ''
				else
					thisCosupervisors = '\\newline\\textit{Co-supervisors: }' .. thisCosupervisors
				end

				local columnOne = thisGraduation

				local thisContents = thisName ..
									'\\newline ' .. thisTitle ..
									'\\newline\\textsc{' .. thisUniversity .. '}, ' .. thisDegree ..
									'\\newline\\textit{Role: }' .. thisRole ..
									thisCosupervisors ..
									'\\newline ' .. thisUrl

				tex.print(columnOne .. ' & ' .. thisContents .. ' \\\\')
			end
			-- Only close the table if one was opened. A group with no entries
			-- never opens a longtable, and closing one that was never opened
			-- aborts the run -- so a CV whose publications happen not to span
			-- every entry type would fail to build at all. (No control
			-- sequences in this comment: luacode scans the body textually for
			-- its terminator, so a stray one here ends the environment early.)
			if not needsHeading then
				tex.print('\\end{longtable}')
				tex.print('\\end{centering}')
			end
		end
	end


	-- Sections come from data/sections.json -- the same file the website reads
	-- -- so the CV cannot quietly omit a section the site is rendering. A
	-- section with no dedicated renderer falls through to prettyprintGeneric,
	-- which is what lets a new section reach the PDF without touching LaTeX.
	local renderers = {
		education   = prettyprintEducation,
		experience  = prettyprintExperience,
		awards      = prettyprintAwards,
		publications= prettyprintPublications,
		conferences = prettyprintConferences,
		software    = prettyprintSoftware,
		teaching    = prettyprintTeaching,
		supervision = prettyprintSupervision,
		generic     = prettyprintGeneric,
	}

	local sections = readJsonOptional('data/sections.json')
	if sections == nil then
		sections = {}
		for _, id in ipairs({'education', 'experience', 'awards', 'publications',
		                     'conferences', 'software', 'teaching', 'supervision'}) do
			table.insert(sections, {id = id})
		end
	end

	for _, section in ipairs(sections) do
		if section.inCV ~= false and section.enable ~= false then
			section.data     = section.data or section.id
			section.style    = section.style or section.id
			section.cvLabel  = section.cvLabel or section.label or firstToUpper(section.id)
			-- A declared section whose data file is absent or empty is skipped
			-- rather than printed as a bare heading with nothing under it. It
			-- also keeps a freshly scaffolded site -- every data file still []
			-- -- building a valid, if short, CV.
			local records = readJsonOptional('data/' .. section.data .. '.json')
			if records ~= nil and #records > 0 then
				local render = renderers[section.style] or prettyprintGeneric
				render(records, section)
			end
		end
	end
\end{luacode}

\end{document}
